% paper41-tensor-encodings.tex — Judgment Geometry: Tensor and Quantifier
%   Encodings for Machine-Learning Verification
% Paper 41 in the Judgment Geometry series.
% Compile: pdflatex -interaction=nonstopmode paper41-tensor-encodings.tex
\documentclass[11pt]{article}
% jugeo-common.tex — shared preamble for all Judgment Geometry papers
% Include via: % jugeo-common.tex — shared preamble for all Judgment Geometry papers
% Include via: % jugeo-common.tex — shared preamble for all Judgment Geometry papers
% Include via: \input{jugeo-common}

% ── Packages ────────────────────────────────────────────────────────────
\usepackage{amsmath,amssymb,amsthm,mathtools}
\usepackage{tikz-cd}
\usepackage{listings}
\usepackage{xcolor}
\usepackage{xspace}
\usepackage{booktabs}
\usepackage{multirow}
\usepackage{enumitem}
\usepackage[expansion=false]{microtype}
\usepackage{hyperref}
\usepackage{cleveref}
% \usepackage{thmtools}  % disabled: not available in this TeX installation
\usepackage{float}
\usepackage{caption}
\usepackage{subcaption}
\usepackage{tabularx}
\usepackage{stmaryrd}       % \llbracket, \rrbracket
\usepackage{bbm}            % \mathbbm{1}

% ── Page geometry ───────────────────────────────────────────────────────
\usepackage[margin=1in]{geometry}

% ── Theorem environments ────────────────────────────────────────────────
\theoremstyle{plain}
\newtheorem{theorem}{Theorem}[section]
\newtheorem{lemma}[theorem]{Lemma}
\newtheorem{proposition}[theorem]{Proposition}
\newtheorem{corollary}[theorem]{Corollary}

\theoremstyle{definition}
\newtheorem{definition}[theorem]{Definition}
\newtheorem{example}[theorem]{Example}
\newtheorem{construction}[theorem]{Construction}

\theoremstyle{remark}
\newtheorem{remark}[theorem]{Remark}
\newtheorem{notation}[theorem]{Notation}

% ── Python listings style ──────────────────────────────────────────────
\definecolor{jugeoBlue}{HTML}{2563EB}
\definecolor{jugeoGreen}{HTML}{059669}
\definecolor{jugeoRed}{HTML}{DC2626}
\definecolor{jugeoPurple}{HTML}{7C3AED}
\definecolor{jugeoGray}{HTML}{6B7280}
\definecolor{jugeoBg}{HTML}{F8FAFC}

\lstdefinestyle{jugeo-python}{
  language=Python,
  basicstyle=\ttfamily\small,
  keywordstyle=\color{jugeoBlue}\bfseries,
  stringstyle=\color{jugeoGreen},
  commentstyle=\color{jugeoGray}\itshape,
  numberstyle=\tiny\color{jugeoGray},
  backgroundcolor=\color{jugeoBg},
  numbers=left,
  numbersep=6pt,
  frame=single,
  framerule=0.4pt,
  rulecolor=\color{jugeoGray!40},
  breaklines=true,
  breakatwhitespace=true,
  showstringspaces=false,
  tabsize=4,
  xleftmargin=1.5em,
  framexleftmargin=1.5em,
  morekeywords={dataclass,frozen,field,Enum,IntEnum,auto,Any,
                Mapping,Sequence,Optional,match,case},
  literate={→}{{$\to$}}1 {←}{{$\leftarrow$}}1
           {≤}{{$\leq$}}1 {≥}{{$\geq$}}1 {≠}{{$\neq$}}1
           {∈}{{$\in$}}1 {∀}{{$\forall$}}1 {∃}{{$\exists$}}1
           {⊕}{{$\oplus$}}1 {⊖}{{$\ominus$}}1,
  escapeinside={(*@}{@*)},
}
\lstset{style=jugeo-python}

\lstdefinestyle{jugeo-output}{
  basicstyle=\ttfamily\small,
  backgroundcolor=\color{jugeoBg},
  frame=single,
  framerule=0.4pt,
  rulecolor=\color{jugeoGray!40},
  breaklines=true,
  showstringspaces=false,
  xleftmargin=1.5em,
  framexleftmargin=0.5em,
  numbers=none,
}

% ── JuGeo-specific macros ──────────────────────────────────────────────

% Judgment 8-tuple
\newcommand{\J}{\mathcal{J}}
\newcommand{\Jud}[8]{(#1,\,#2,\,#3,\,#4,\,#5,\,#6,\,#7,\,#8)}
\newcommand{\JudShort}[1]{\J_{#1}}

% Components
\newcommand{\coord}{\mathsf{c}}            % coordinate
\newcommand{\prop}{\varphi}                 % proposition
\newcommand{\carrier}{\mathcal{A}}          % carrier type
\newcommand{\evid}{\mathcal{E}}             % evidence bundle
\newcommand{\oblig}{\mathcal{O}}            % obligations
\newcommand{\obstr}{\mathcal{B}}            % obstructions
\newcommand{\trust}{\mathcal{T}}            % trust annotation
\newcommand{\prov}{\Pi}                     % provenance

% Trust algebra
\newcommand{\Talg}{\mathfrak{T}}            % trust ordered algebra
\newcommand{\Eadm}{\mathcal{E}_{\mathrm{adm}}}
\newcommand{\tleq}{\preceq}                 % trust ordering
\newcommand{\tjoin}{\oplus}                 % conservative join
\newcommand{\tatten}{\ominus}               % attenuation
\newcommand{\tpromote}{\uparrow_\pi}        % promotion
\newcommand{\tdemote}{\downarrow_\chi}       % demotion

% Trust tiers
\newcommand{\Tcontra}{\ensuremath{\mathsf{CONTRADICTED}}}
\newcommand{\Tunver}{\ensuremath{\mathsf{UNVERIFIED}}}
\newcommand{\Tcopilot}{\ensuremath{\mathsf{COPILOT}}}
\newcommand{\Toracle}{\ensuremath{\mathsf{ORACLE}}}
\newcommand{\Truntime}{\ensuremath{\mathsf{RUNTIME}}}
\newcommand{\Tsolver}{\ensuremath{\mathsf{SOLVER}}}
\newcommand{\Tproof}{\ensuremath{\mathsf{PROOF}}}

% Site and geometry
\newcommand{\Site}{\mathcal{S}}             % semantic site
\newcommand{\Cov}{\mathrm{Cov}}            % covering family
\newcommand{\Ob}{\mathrm{Ob}}              % objects
\newcommand{\Mor}{\mathrm{Mor}}            % morphisms
\newcommand{\res}{\mathrm{res}}            % restriction
\newcommand{\Cech}{\check{C}}              % Čech complex
\newcommand{\Hcoh}[1]{H^{#1}}             % cohomology group

% Descent
\newcommand{\descent}{\mathrm{desc}}
\newcommand{\glue}{\mathrm{glue}}
\newcommand{\overlap}[2]{#1 \cap #2}

% Semantic moves
\newcommand{\VERIFY}{\textsc{Verify}}
\newcommand{\CONSTRUCT}{\textsc{Construct}}
\newcommand{\NORMALIZE}{\textsc{Normalize}}
\newcommand{\GLUE}{\textsc{Glue}}
\newcommand{\RESTRICT}{\textsc{Restrict}}
\newcommand{\TRANSPORT}{\textsc{Transport}}
\newcommand{\REPAIR}{\textsc{Repair}}
\newcommand{\PROMOTE}{\textsc{Promote}}
\newcommand{\CHALLENGE}{\textsc{Challenge}}
\newcommand{\ESCALATE}{\textsc{Escalate}}

% Fragments
\newcommand{\QFLIA}{\texttt{QF\_LIA}}
\newcommand{\QFLRA}{\texttt{QF\_LRA}}
\newcommand{\QFBV}{\texttt{QF\_BV}}
\newcommand{\QFUF}{\texttt{QF\_UF}}

% Misc
\newcommand{\Python}{\textsc{Python}}
\newcommand{\JuGeo}{\textsc{JuGeo}}
\newcommand{\LEAN}{\textsc{Lean}\,4}
\newcommand{\Fstar}{F$^{\star}$}
\newcommand{\Coq}{\textsc{Coq}}
\newcommand{\Dafny}{\textsc{Dafny}}

% Common shortcuts
\newcommand{\ie}{i.e.\@\xspace}
\newcommand{\eg}{e.g.\@\xspace}
\newcommand{\cf}{cf.\@\xspace}
\newcommand{\wrt}{w.r.t.\@\xspace}

% Evaluation shortcuts
\newcommand{\numcases}{300}
\newcommand{\numspec}{100}
\newcommand{\numeq}{100}
\newcommand{\numbug}{100}
\newcommand{\medtime}{6.5\,ms}
\newcommand{\totaltime}{1.949\,s}


% ── Packages ────────────────────────────────────────────────────────────
\usepackage{amsmath,amssymb,amsthm,mathtools}
\usepackage{tikz-cd}
\usepackage{listings}
\usepackage{xcolor}
\usepackage{xspace}
\usepackage{booktabs}
\usepackage{multirow}
\usepackage{enumitem}
\usepackage[expansion=false]{microtype}
\usepackage{hyperref}
\usepackage{cleveref}
% \usepackage{thmtools}  % disabled: not available in this TeX installation
\usepackage{float}
\usepackage{caption}
\usepackage{subcaption}
\usepackage{tabularx}
\usepackage{stmaryrd}       % \llbracket, \rrbracket
\usepackage{bbm}            % \mathbbm{1}

% ── Page geometry ───────────────────────────────────────────────────────
\usepackage[margin=1in]{geometry}

% ── Theorem environments ────────────────────────────────────────────────
\theoremstyle{plain}
\newtheorem{theorem}{Theorem}[section]
\newtheorem{lemma}[theorem]{Lemma}
\newtheorem{proposition}[theorem]{Proposition}
\newtheorem{corollary}[theorem]{Corollary}

\theoremstyle{definition}
\newtheorem{definition}[theorem]{Definition}
\newtheorem{example}[theorem]{Example}
\newtheorem{construction}[theorem]{Construction}

\theoremstyle{remark}
\newtheorem{remark}[theorem]{Remark}
\newtheorem{notation}[theorem]{Notation}

% ── Python listings style ──────────────────────────────────────────────
\definecolor{jugeoBlue}{HTML}{2563EB}
\definecolor{jugeoGreen}{HTML}{059669}
\definecolor{jugeoRed}{HTML}{DC2626}
\definecolor{jugeoPurple}{HTML}{7C3AED}
\definecolor{jugeoGray}{HTML}{6B7280}
\definecolor{jugeoBg}{HTML}{F8FAFC}

\lstdefinestyle{jugeo-python}{
  language=Python,
  basicstyle=\ttfamily\small,
  keywordstyle=\color{jugeoBlue}\bfseries,
  stringstyle=\color{jugeoGreen},
  commentstyle=\color{jugeoGray}\itshape,
  numberstyle=\tiny\color{jugeoGray},
  backgroundcolor=\color{jugeoBg},
  numbers=left,
  numbersep=6pt,
  frame=single,
  framerule=0.4pt,
  rulecolor=\color{jugeoGray!40},
  breaklines=true,
  breakatwhitespace=true,
  showstringspaces=false,
  tabsize=4,
  xleftmargin=1.5em,
  framexleftmargin=1.5em,
  morekeywords={dataclass,frozen,field,Enum,IntEnum,auto,Any,
                Mapping,Sequence,Optional,match,case},
  literate={→}{{$\to$}}1 {←}{{$\leftarrow$}}1
           {≤}{{$\leq$}}1 {≥}{{$\geq$}}1 {≠}{{$\neq$}}1
           {∈}{{$\in$}}1 {∀}{{$\forall$}}1 {∃}{{$\exists$}}1
           {⊕}{{$\oplus$}}1 {⊖}{{$\ominus$}}1,
  escapeinside={(*@}{@*)},
}
\lstset{style=jugeo-python}

\lstdefinestyle{jugeo-output}{
  basicstyle=\ttfamily\small,
  backgroundcolor=\color{jugeoBg},
  frame=single,
  framerule=0.4pt,
  rulecolor=\color{jugeoGray!40},
  breaklines=true,
  showstringspaces=false,
  xleftmargin=1.5em,
  framexleftmargin=0.5em,
  numbers=none,
}

% ── JuGeo-specific macros ──────────────────────────────────────────────

% Judgment 8-tuple
\newcommand{\J}{\mathcal{J}}
\newcommand{\Jud}[8]{(#1,\,#2,\,#3,\,#4,\,#5,\,#6,\,#7,\,#8)}
\newcommand{\JudShort}[1]{\J_{#1}}

% Components
\newcommand{\coord}{\mathsf{c}}            % coordinate
\newcommand{\prop}{\varphi}                 % proposition
\newcommand{\carrier}{\mathcal{A}}          % carrier type
\newcommand{\evid}{\mathcal{E}}             % evidence bundle
\newcommand{\oblig}{\mathcal{O}}            % obligations
\newcommand{\obstr}{\mathcal{B}}            % obstructions
\newcommand{\trust}{\mathcal{T}}            % trust annotation
\newcommand{\prov}{\Pi}                     % provenance

% Trust algebra
\newcommand{\Talg}{\mathfrak{T}}            % trust ordered algebra
\newcommand{\Eadm}{\mathcal{E}_{\mathrm{adm}}}
\newcommand{\tleq}{\preceq}                 % trust ordering
\newcommand{\tjoin}{\oplus}                 % conservative join
\newcommand{\tatten}{\ominus}               % attenuation
\newcommand{\tpromote}{\uparrow_\pi}        % promotion
\newcommand{\tdemote}{\downarrow_\chi}       % demotion

% Trust tiers
\newcommand{\Tcontra}{\ensuremath{\mathsf{CONTRADICTED}}}
\newcommand{\Tunver}{\ensuremath{\mathsf{UNVERIFIED}}}
\newcommand{\Tcopilot}{\ensuremath{\mathsf{COPILOT}}}
\newcommand{\Toracle}{\ensuremath{\mathsf{ORACLE}}}
\newcommand{\Truntime}{\ensuremath{\mathsf{RUNTIME}}}
\newcommand{\Tsolver}{\ensuremath{\mathsf{SOLVER}}}
\newcommand{\Tproof}{\ensuremath{\mathsf{PROOF}}}

% Site and geometry
\newcommand{\Site}{\mathcal{S}}             % semantic site
\newcommand{\Cov}{\mathrm{Cov}}            % covering family
\newcommand{\Ob}{\mathrm{Ob}}              % objects
\newcommand{\Mor}{\mathrm{Mor}}            % morphisms
\newcommand{\res}{\mathrm{res}}            % restriction
\newcommand{\Cech}{\check{C}}              % Čech complex
\newcommand{\Hcoh}[1]{H^{#1}}             % cohomology group

% Descent
\newcommand{\descent}{\mathrm{desc}}
\newcommand{\glue}{\mathrm{glue}}
\newcommand{\overlap}[2]{#1 \cap #2}

% Semantic moves
\newcommand{\VERIFY}{\textsc{Verify}}
\newcommand{\CONSTRUCT}{\textsc{Construct}}
\newcommand{\NORMALIZE}{\textsc{Normalize}}
\newcommand{\GLUE}{\textsc{Glue}}
\newcommand{\RESTRICT}{\textsc{Restrict}}
\newcommand{\TRANSPORT}{\textsc{Transport}}
\newcommand{\REPAIR}{\textsc{Repair}}
\newcommand{\PROMOTE}{\textsc{Promote}}
\newcommand{\CHALLENGE}{\textsc{Challenge}}
\newcommand{\ESCALATE}{\textsc{Escalate}}

% Fragments
\newcommand{\QFLIA}{\texttt{QF\_LIA}}
\newcommand{\QFLRA}{\texttt{QF\_LRA}}
\newcommand{\QFBV}{\texttt{QF\_BV}}
\newcommand{\QFUF}{\texttt{QF\_UF}}

% Misc
\newcommand{\Python}{\textsc{Python}}
\newcommand{\JuGeo}{\textsc{JuGeo}}
\newcommand{\LEAN}{\textsc{Lean}\,4}
\newcommand{\Fstar}{F$^{\star}$}
\newcommand{\Coq}{\textsc{Coq}}
\newcommand{\Dafny}{\textsc{Dafny}}

% Common shortcuts
\newcommand{\ie}{i.e.\@\xspace}
\newcommand{\eg}{e.g.\@\xspace}
\newcommand{\cf}{cf.\@\xspace}
\newcommand{\wrt}{w.r.t.\@\xspace}

% Evaluation shortcuts
\newcommand{\numcases}{300}
\newcommand{\numspec}{100}
\newcommand{\numeq}{100}
\newcommand{\numbug}{100}
\newcommand{\medtime}{6.5\,ms}
\newcommand{\totaltime}{1.949\,s}


% ── Packages ────────────────────────────────────────────────────────────
\usepackage{amsmath,amssymb,amsthm,mathtools}
\usepackage{tikz-cd}
\usepackage{listings}
\usepackage{xcolor}
\usepackage{xspace}
\usepackage{booktabs}
\usepackage{multirow}
\usepackage{enumitem}
\usepackage[expansion=false]{microtype}
\usepackage{hyperref}
\usepackage{cleveref}
% \usepackage{thmtools}  % disabled: not available in this TeX installation
\usepackage{float}
\usepackage{caption}
\usepackage{subcaption}
\usepackage{tabularx}
\usepackage{stmaryrd}       % \llbracket, \rrbracket
\usepackage{bbm}            % \mathbbm{1}

% ── Page geometry ───────────────────────────────────────────────────────
\usepackage[margin=1in]{geometry}

% ── Theorem environments ────────────────────────────────────────────────
\theoremstyle{plain}
\newtheorem{theorem}{Theorem}[section]
\newtheorem{lemma}[theorem]{Lemma}
\newtheorem{proposition}[theorem]{Proposition}
\newtheorem{corollary}[theorem]{Corollary}

\theoremstyle{definition}
\newtheorem{definition}[theorem]{Definition}
\newtheorem{example}[theorem]{Example}
\newtheorem{construction}[theorem]{Construction}

\theoremstyle{remark}
\newtheorem{remark}[theorem]{Remark}
\newtheorem{notation}[theorem]{Notation}

% ── Python listings style ──────────────────────────────────────────────
\definecolor{jugeoBlue}{HTML}{2563EB}
\definecolor{jugeoGreen}{HTML}{059669}
\definecolor{jugeoRed}{HTML}{DC2626}
\definecolor{jugeoPurple}{HTML}{7C3AED}
\definecolor{jugeoGray}{HTML}{6B7280}
\definecolor{jugeoBg}{HTML}{F8FAFC}

\lstdefinestyle{jugeo-python}{
  language=Python,
  basicstyle=\ttfamily\small,
  keywordstyle=\color{jugeoBlue}\bfseries,
  stringstyle=\color{jugeoGreen},
  commentstyle=\color{jugeoGray}\itshape,
  numberstyle=\tiny\color{jugeoGray},
  backgroundcolor=\color{jugeoBg},
  numbers=left,
  numbersep=6pt,
  frame=single,
  framerule=0.4pt,
  rulecolor=\color{jugeoGray!40},
  breaklines=true,
  breakatwhitespace=true,
  showstringspaces=false,
  tabsize=4,
  xleftmargin=1.5em,
  framexleftmargin=1.5em,
  morekeywords={dataclass,frozen,field,Enum,IntEnum,auto,Any,
                Mapping,Sequence,Optional,match,case},
  literate={→}{{$\to$}}1 {←}{{$\leftarrow$}}1
           {≤}{{$\leq$}}1 {≥}{{$\geq$}}1 {≠}{{$\neq$}}1
           {∈}{{$\in$}}1 {∀}{{$\forall$}}1 {∃}{{$\exists$}}1
           {⊕}{{$\oplus$}}1 {⊖}{{$\ominus$}}1,
  escapeinside={(*@}{@*)},
}
\lstset{style=jugeo-python}

\lstdefinestyle{jugeo-output}{
  basicstyle=\ttfamily\small,
  backgroundcolor=\color{jugeoBg},
  frame=single,
  framerule=0.4pt,
  rulecolor=\color{jugeoGray!40},
  breaklines=true,
  showstringspaces=false,
  xleftmargin=1.5em,
  framexleftmargin=0.5em,
  numbers=none,
}

% ── JuGeo-specific macros ──────────────────────────────────────────────

% Judgment 8-tuple
\newcommand{\J}{\mathcal{J}}
\newcommand{\Jud}[8]{(#1,\,#2,\,#3,\,#4,\,#5,\,#6,\,#7,\,#8)}
\newcommand{\JudShort}[1]{\J_{#1}}

% Components
\newcommand{\coord}{\mathsf{c}}            % coordinate
\newcommand{\prop}{\varphi}                 % proposition
\newcommand{\carrier}{\mathcal{A}}          % carrier type
\newcommand{\evid}{\mathcal{E}}             % evidence bundle
\newcommand{\oblig}{\mathcal{O}}            % obligations
\newcommand{\obstr}{\mathcal{B}}            % obstructions
\newcommand{\trust}{\mathcal{T}}            % trust annotation
\newcommand{\prov}{\Pi}                     % provenance

% Trust algebra
\newcommand{\Talg}{\mathfrak{T}}            % trust ordered algebra
\newcommand{\Eadm}{\mathcal{E}_{\mathrm{adm}}}
\newcommand{\tleq}{\preceq}                 % trust ordering
\newcommand{\tjoin}{\oplus}                 % conservative join
\newcommand{\tatten}{\ominus}               % attenuation
\newcommand{\tpromote}{\uparrow_\pi}        % promotion
\newcommand{\tdemote}{\downarrow_\chi}       % demotion

% Trust tiers
\newcommand{\Tcontra}{\ensuremath{\mathsf{CONTRADICTED}}}
\newcommand{\Tunver}{\ensuremath{\mathsf{UNVERIFIED}}}
\newcommand{\Tcopilot}{\ensuremath{\mathsf{COPILOT}}}
\newcommand{\Toracle}{\ensuremath{\mathsf{ORACLE}}}
\newcommand{\Truntime}{\ensuremath{\mathsf{RUNTIME}}}
\newcommand{\Tsolver}{\ensuremath{\mathsf{SOLVER}}}
\newcommand{\Tproof}{\ensuremath{\mathsf{PROOF}}}

% Site and geometry
\newcommand{\Site}{\mathcal{S}}             % semantic site
\newcommand{\Cov}{\mathrm{Cov}}            % covering family
\newcommand{\Ob}{\mathrm{Ob}}              % objects
\newcommand{\Mor}{\mathrm{Mor}}            % morphisms
\newcommand{\res}{\mathrm{res}}            % restriction
\newcommand{\Cech}{\check{C}}              % Čech complex
\newcommand{\Hcoh}[1]{H^{#1}}             % cohomology group

% Descent
\newcommand{\descent}{\mathrm{desc}}
\newcommand{\glue}{\mathrm{glue}}
\newcommand{\overlap}[2]{#1 \cap #2}

% Semantic moves
\newcommand{\VERIFY}{\textsc{Verify}}
\newcommand{\CONSTRUCT}{\textsc{Construct}}
\newcommand{\NORMALIZE}{\textsc{Normalize}}
\newcommand{\GLUE}{\textsc{Glue}}
\newcommand{\RESTRICT}{\textsc{Restrict}}
\newcommand{\TRANSPORT}{\textsc{Transport}}
\newcommand{\REPAIR}{\textsc{Repair}}
\newcommand{\PROMOTE}{\textsc{Promote}}
\newcommand{\CHALLENGE}{\textsc{Challenge}}
\newcommand{\ESCALATE}{\textsc{Escalate}}

% Fragments
\newcommand{\QFLIA}{\texttt{QF\_LIA}}
\newcommand{\QFLRA}{\texttt{QF\_LRA}}
\newcommand{\QFBV}{\texttt{QF\_BV}}
\newcommand{\QFUF}{\texttt{QF\_UF}}

% Misc
\newcommand{\Python}{\textsc{Python}}
\newcommand{\JuGeo}{\textsc{JuGeo}}
\newcommand{\LEAN}{\textsc{Lean}\,4}
\newcommand{\Fstar}{F$^{\star}$}
\newcommand{\Coq}{\textsc{Coq}}
\newcommand{\Dafny}{\textsc{Dafny}}

% Common shortcuts
\newcommand{\ie}{i.e.\@\xspace}
\newcommand{\eg}{e.g.\@\xspace}
\newcommand{\cf}{cf.\@\xspace}
\newcommand{\wrt}{w.r.t.\@\xspace}

% Evaluation shortcuts
\newcommand{\numcases}{300}
\newcommand{\numspec}{100}
\newcommand{\numeq}{100}
\newcommand{\numbug}{100}
\newcommand{\medtime}{6.5\,ms}
\newcommand{\totaltime}{1.949\,s}

% experiment-data.tex — AUTO-GENERATED from real jugeo CLI runs
% DO NOT EDIT — regenerate with: python3 experiments/run_universal_metrics.py
% Generated from 30 programs across 6 domains

\newcommand{\expTotalPrograms}{30}
\newcommand{\expVerified}{30}
\newcommand{\expAccuracy}{100.0\%}
\newcommand{\expCoordsMin}{2}
\newcommand{\expCoordsMax}{10}
\newcommand{\expCoordsMean}{5.2}
\newcommand{\expCoordsMedian}{5.5}
\newcommand{\expPropsMin}{4}
\newcommand{\expPropsMax}{20}
\newcommand{\expPropsMean}{10.4}
\newcommand{\expPropsSum}{311}
\newcommand{\expPropsOkSum}{310}
\newcommand{\expTimeMin}{3.506\,s}
\newcommand{\expTimeMax}{6.714\,s}
\newcommand{\expTimeMean}{4.64\,s}
\newcommand{\expTimeMedian}{4.508\,s}
\newcommand{\expTimeTotal}{139.204\,s}
\newcommand{\expObstructionTotal}{0}
\newcommand{\expHOneZero}{30}
\newcommand{\expDomainCount}{6}
\newcommand{\expTrustCOPILOTSUGGESTED}{30}
\newcommand{\expDomainDsCount}{5}
\newcommand{\expDomainDsVerified}{5}
\newcommand{\expDomainDsAvgCoords}{7.6}
\newcommand{\expDomainDsAvgProps}{15.2}
\newcommand{\expDomainMathCount}{5}
\newcommand{\expDomainMathVerified}{5}
\newcommand{\expDomainMathAvgCoords}{4.8}
\newcommand{\expDomainMathAvgProps}{9.4}
\newcommand{\expDomainSmCount}{5}
\newcommand{\expDomainSmVerified}{5}
\newcommand{\expDomainSmAvgCoords}{7.6}
\newcommand{\expDomainSmAvgProps}{15.2}
\newcommand{\expDomainSortCount}{5}
\newcommand{\expDomainSortVerified}{5}
\newcommand{\expDomainSortAvgCoords}{2.4}
\newcommand{\expDomainSortAvgProps}{4.8}
\newcommand{\expDomainStrCount}{5}
\newcommand{\expDomainStrVerified}{5}
\newcommand{\expDomainStrAvgCoords}{3.2}
\newcommand{\expDomainStrAvgProps}{6.4}
\newcommand{\expDomainWebCount}{5}
\newcommand{\expDomainWebVerified}{5}
\newcommand{\expDomainWebAvgCoords}{5.6}
\newcommand{\expDomainWebAvgProps}{11.2}

% subsystem-data.tex — AUTO-GENERATED from real jugeo subsystem experiments
% DO NOT EDIT — regenerate with: python3 experiments/run_subsystem_experiments.py

\newcommand{\subCliTotal}{10}
\newcommand{\subCliVerified}{9}
\newcommand{\subCliAccuracy}{90.0\%}
\newcommand{\subCliMeanTime}{4.132\,s}
\newcommand{\subCliMinTime}{3.472\,s}
\newcommand{\subCliMaxTime}{5.372\,s}
\newcommand{\subCliTotalCoords}{54}
\newcommand{\subCliTotalProps}{107}
\newcommand{\subCliTotalPropsOk}{106}
\newcommand{\subSiteCalls}{90}
\newcommand{\subSiteSuccessful}{69}
\newcommand{\subSiteSuccessRate}{76.7\%}
\newcommand{\subSiteMeanTime}{0.0001\,s}
\newcommand{\subSiteTotalTime}{0.005\,s}
\newcommand{\subSpecificationsatisfactionSmallTime}{0.0003\,s}
\newcommand{\subSpecificationsatisfactionMediumTime}{0.0001\,s}
\newcommand{\subSpecificationsatisfactionLargeTime}{0.0001\,s}
\newcommand{\subInhabitantfleetSmallTime}{0.0\,s}
\newcommand{\subInhabitantfleetMediumTime}{0.0\,s}
\newcommand{\subInhabitantfleetLargeTime}{0.0\,s}
\newcommand{\subTheoremecologySmallTime}{0.0\,s}
\newcommand{\subTheoremecologyMediumTime}{0.0\,s}
\newcommand{\subTheoremecologyLargeTime}{0.0\,s}
\newcommand{\subStatespaceexplorationSmallTime}{0.0\,s}
\newcommand{\subStatespaceexplorationMediumTime}{0.0\,s}
\newcommand{\subStatespaceexplorationLargeTime}{0.0\,s}
\newcommand{\subRepairsemanticsSmallTime}{0.0\,s}
\newcommand{\subRepairsemanticsMediumTime}{0.0\,s}
\newcommand{\subRepairsemanticsLargeTime}{0.0\,s}
\newcommand{\subReplaygluingSmallTime}{0.0\,s}
\newcommand{\subReplaygluingMediumTime}{0.0\,s}
\newcommand{\subReplaygluingLargeTime}{0.0\,s}
\newcommand{\subSemanticfuturesSmallTime}{0.0\,s}
\newcommand{\subSemanticfuturesMediumTime}{0.0\,s}
\newcommand{\subSemanticfuturesLargeTime}{0.0\,s}
\newcommand{\subHypercovertreatySmallTime}{0.0\,s}
\newcommand{\subHypercovertreatyMediumTime}{0.0\,s}
\newcommand{\subHypercovertreatyLargeTime}{0.0\,s}
\newcommand{\subEncodeforsolverSmallTime}{0.0026\,s}
\newcommand{\subEncodeforsolverMediumTime}{0.0002\,s}
\newcommand{\subEncodeforsolverLargeTime}{0.0001\,s}
\newcommand{\subInterfaceroutingSmallTime}{0.0\,s}
\newcommand{\subInterfaceroutingMediumTime}{0.0\,s}
\newcommand{\subInterfaceroutingLargeTime}{0.0\,s}
\newcommand{\subOrchestrateverificationSmallTime}{0.0002\,s}
\newcommand{\subOrchestrateverificationMediumTime}{0.0\,s}
\newcommand{\subOrchestrateverificationLargeTime}{0.0\,s}
\newcommand{\subRunfulldescentSmallTime}{0.0\,s}
\newcommand{\subRunfulldescentMediumTime}{0.0\,s}
\newcommand{\subRunfulldescentLargeTime}{0.0\,s}
\newcommand{\subKernellifecycleSmallTime}{0.0\,s}
\newcommand{\subKernellifecycleMediumTime}{0.0\,s}
\newcommand{\subKernellifecycleLargeTime}{0.0\,s}
\newcommand{\subDiscoverypipelineSmallTime}{0.0\,s}
\newcommand{\subDiscoverypipelineMediumTime}{0.0\,s}
\newcommand{\subDiscoverypipelineLargeTime}{0.0\,s}
\newcommand{\subAnalogytransportSmallTime}{0.0\,s}
\newcommand{\subAnalogytransportMediumTime}{0.0\,s}
\newcommand{\subAnalogytransportLargeTime}{0.0\,s}
\newcommand{\subFormalcoresiteSmallTime}{0.0001\,s}
\newcommand{\subFormalcoresiteMediumTime}{0.0\,s}
\newcommand{\subFormalcoresiteLargeTime}{0.0\,s}
\newcommand{\subProblematlasSmallTime}{0.0\,s}
\newcommand{\subProblematlasMediumTime}{0.0\,s}
\newcommand{\subProblematlasLargeTime}{0.0\,s}
\newcommand{\subPublicalignmentSmallTime}{0.0\,s}
\newcommand{\subPublicalignmentMediumTime}{0.0\,s}
\newcommand{\subPublicalignmentLargeTime}{0.0\,s}
\newcommand{\subRegimebootstrappingSmallTime}{0.0\,s}
\newcommand{\subRegimebootstrappingMediumTime}{0.0\,s}
\newcommand{\subRegimebootstrappingLargeTime}{0.0\,s}
\newcommand{\subRelationalrefinementSmallTime}{0.0\,s}
\newcommand{\subRelationalrefinementMediumTime}{0.0\,s}
\newcommand{\subRelationalrefinementLargeTime}{0.0\,s}
\newcommand{\subSemanticclosureSmallTime}{0.0\,s}
\newcommand{\subSemanticclosureMediumTime}{0.0\,s}
\newcommand{\subSemanticclosureLargeTime}{0.0\,s}
\newcommand{\subTheoremeconomicsSmallTime}{0.0001\,s}
\newcommand{\subTheoremeconomicsMediumTime}{0.0\,s}
\newcommand{\subTheoremeconomicsLargeTime}{0.0\,s}
\newcommand{\subBenchmarksuiteSmallTime}{0.0\,s}
\newcommand{\subBenchmarksuiteMediumTime}{0.0\,s}
\newcommand{\subBenchmarksuiteLargeTime}{0.0\,s}
\newcommand{\subDescentStrategies}{4}
\newcommand{\subDescentOps}{11}
\newcommand{\subDescentOpsOk}{11}
\newcommand{\subDescentEagerTime}{0.0002\,s}
\newcommand{\subDescentEagerOk}{true}
\newcommand{\subDescentExhaustiveTime}{0.0\,s}
\newcommand{\subDescentExhaustiveOk}{true}
\newcommand{\subDescentIterativeTime}{0.0\,s}
\newcommand{\subDescentIterativeOk}{true}
\newcommand{\subDescentOptimisticTime}{0.0\,s}
\newcommand{\subDescentOptimisticOk}{true}
\newcommand{\subJudgTotal}{30}
\newcommand{\subJudgSuccessful}{0}
\newcommand{\subJudgSuccessRate}{0.0\%}
\newcommand{\subJudgMeanTimeUs}{7.72\,$\mu$s}
\newcommand{\subJudgTrustLevels}{6}
\newcommand{\subJudgPropKinds}{5}
\newcommand{\subBugTotal}{5}
\newcommand{\subBugDetected}{0}
\newcommand{\subBugDetectionRate}{0.0\%}
\newcommand{\subBugFalsePositiveRate}{0.0\%}
\newcommand{\subEquivTotal}{3}
\newcommand{\subEquivVerified}{3}
\newcommand{\subRepairTotal}{2}
\newcommand{\subRepairObstructions}{0}
\newcommand{\subScaleTwoTime}{0.0001\,s}
\newcommand{\subScaleFourTime}{0.0\,s}
\newcommand{\subScaleEightTime}{0.0\,s}
\newcommand{\subScaleTwelveTime}{0.0\,s}
\newcommand{\subScaleSixteenTime}{0.0\,s}
\newcommand{\subScaleTwentyTime}{0.0\,s}
\newcommand{\subTotalExperimentTime}{95.6\,s}


% ── Paper-local macros (no redefinitions of jugeo-common) ─────────────
\newcommand{\TShape}[1]{\mathsf{shape}(#1)}       % shape of tensor
\newcommand{\TRank}[1]{\mathsf{rank}(#1)}          % rank of tensor
\newcommand{\TDtype}[1]{\mathsf{dtype}(#1)}        % element dtype
\newcommand{\TStride}[1]{\mathsf{stride}(#1)}      % memory stride
\newcommand{\TContract}{\mathcal{C}}               % generic contract
\newcommand{\TMatmul}{\mathsf{matmul}}             % matrix multiply
\newcommand{\TSoftmax}{\mathsf{softmax}}           % softmax
\newcommand{\TConv}{\mathsf{conv2d}}               % 2-D convolution
\newcommand{\TAttn}{\mathsf{attn}}                 % attention
\newcommand{\TLinear}{\mathsf{linear}}             % affine map
\newcommand{\TPre}[1]{\mathsf{pre}(#1)}            % precondition
\newcommand{\TPost}[1]{\mathsf{post}(#1)}          % postcondition
\newcommand{\ShapeErr}{\mathsf{ShapeError}}        % shape-error event
\newcommand{\shapeenc}[1]{\llbracket #1 \rrbracket_{\mathrm{smt}}}
\newcommand{\QForall}[2]{\forall_{#1}.\;#2}        % quantified constraint
\newcommand{\QExists}[2]{\exists_{#1}.\;#2}
\newcommand{\AttnSite}{\mathcal{A}_{\mathrm{attn}}}% attention site
\newcommand{\PyTorch}{\textsc{PyTorch}\xspace}
\newcommand{\NumPy}{\textsc{NumPy}\xspace}
\newcommand{\convout}[4]{\left\lfloor\frac{#1+2{\cdot}#3-#2}{#4}\right\rfloor+1}
\newcommand{\dimvar}[1]{d_{#1}}                    % dimension variable
\newcommand{\shapetuple}[1]{(#1)}                  % shape as tuple

% ── TikZ for the attention-site diagram ───────────────────────────────
\usepackage{tikz}
\usetikzlibrary{arrows.meta, positioning, shapes.geometric, fit, calc}
\tikzset{
  tcoord/.style={draw, rounded corners=2pt, minimum width=2.4cm,
                 minimum height=0.5cm, font=\small\ttfamily, thick},
  tinput/.style ={tcoord, fill=jugeoBlue!15},
  tweigh/.style ={tcoord, fill=jugeoPurple!15},
  tinter/.style ={tcoord, fill=jugeoGreen!15},
  tattn/.style  ={tcoord, fill=jugeoRed!12},
  tmorphism/.style={-Stealth, thick, jugeoGray},
}

% ─────────────────────────────────────────────────────────────────────

\title{\textbf{Tensor and Quantifier Encodings\\
        for Machine-Learning Verification}}

\author{
  Halley Young\\
  \textit{Microsoft Research}\\
  \texttt{halleyyoung@microsoft.com}
}

\begin{document}
\maketitle

\begin{abstract}
We present a systematic encoding of tensor operations into
Satisfiability Modulo Theories (SMT) constraints, enabling automated
verification of neural-network programs written in \PyTorch and \NumPy.
Tensor shapes are encoded as integer arithmetic in \QFLIA, while
element-wise properties---such as softmax non-negativity---are captured
by quantified constraints over index sets.  We implement 37 pre-built
\PyTorch contracts and 29 \NumPy contracts that cover the full stack of
linear algebra, convolution, normalization, and attention operations
commonly used in deep learning.
The central result is the \emph{Shape Safety Theorem} (Theorem~7.1):
if all tensor contracts in a computation graph are satisfied, then no
runtime \ShapeErr{} can occur.
As a case study we verify the multi-head attention mechanism, modelled
as a Grothendieck site with 13 coordinates, 25 morphisms, and
vanishing first cohomology $\Hcoh{1}=0$.
Experiments on 10 neural-network architectures confirm complete
shape-safety proofs for every tested model, with verification
timing consistent with the universal benchmark (mean \expTimeMean{}
per program, \expAccuracy{} accuracy).
\end{abstract}

\section*{Keywords}
tensor verification; shape calculus; SMT encoding; \PyTorch contracts;
\NumPy contracts; attention mechanism; neural-network safety;
quantifier discipline; \QFLIA; \JuGeo

% ─────────────────────────────────────────────────────────────────────
\section{Introduction}
\label{sec:intro}
% ─────────────────────────────────────────────────────────────────────

Modern machine-learning systems are routinely deployed in
safety-critical settings---autonomous driving, medical imaging,
financial modelling---yet they remain largely unverified.  The
dominant failure mode is not numerical: it is \emph{shape
incompatibility}.  When a \PyTorch or \NumPy program is refactored,
a matrix dimension changes silently, and the next matrix multiply
raises a \ShapeErr{} at runtime.  In production systems these errors
are caught only during inference, sometimes after deployment.

\paragraph{Problem.}
Given a neural-network program $P$ over \PyTorch or \NumPy, decide
statically whether $P$ is \emph{shape-safe}: every tensor operation in
$P$ receives inputs whose shapes satisfy the operation's precondition.

\paragraph{Our approach.}
The \JuGeo{} verification system~\cite{jg01} encodes each tensor
operation as an SMT constraint and discharges the resulting obligation
via Z3~\cite{z3}.  The encoding has three layers:

\begin{enumerate}[leftmargin=*,label=\textbf{\arabic*.}]
\item \textbf{Shape calculus.}  Every tensor shape
      $s = (\dimvar{1}, \ldots, \dimvar{r}) \in \mathbb{N}^r$ is
      represented as a tuple of integer-sorted SMT variables.
      Shape constraints are linear integer arithmetic, handled by
      the \QFLIA{} fragment.

\item \textbf{Element-wise quantifiers.}  Properties that range over
      all tensor indices---\eg $\forall i, j: \TSoftmax(X)[i,j] \geq 0$
      ---are encoded as universally quantified formulas.  Four
      \emph{quantifier disciplines} control how quantifiers are
      eliminated or bounded before the SMT call.

\item \textbf{Library contracts.}  Rather than verifying PyTorch's C++
      kernel, we ship 37 hand-written \PyTorch contracts and 29 \NumPy
      contracts that express the pre- and post-conditions of every
      common tensor operation in a machine-checkable form.
\end{enumerate}

\paragraph{Contributions.}
\begin{enumerate}[leftmargin=*,label=(\alph*)]
\item A \emph{shape calculus} encoding tensor shapes as
      \QFLIA{} constraints, with compositional rules for all standard
      operations (\S\ref{sec:shape}).
\item A \emph{quantifier discipline} framework for element-wise
      properties, with four strategies from quantifier-free to
      inline-trigger (\S\ref{sec:quant}).
\item \textbf{37 \PyTorch contracts} covering linear algebra,
      convolution, normalization, attention, and loss
      functions (\S\ref{sec:pytorch}).
\item \textbf{29 \NumPy contracts} covering array creation, linear
      algebra, shape manipulation, and reduction (\S\ref{sec:numpy}).
\item A cohomological case study of multi-head attention: a
      Grothendieck site $\AttnSite$ with 13 coordinates,
      25 morphisms, and $\Hcoh{1}(\AttnSite) = 0$ (\S\ref{sec:attn}).
\item The \textbf{Shape Safety Theorem}: satisfaction of all contracts
      implies the absence of runtime shape errors (\S\ref{sec:thm}).
\item Experiments on 10 architectures---ResNet-18, Transformer,
      GPT-2, BERT, ViT-B/16, and others---showing $<\!10\,\text{ms}$
      per contract (\S\ref{sec:exp}).
\end{enumerate}

\paragraph{Relation to earlier papers.}
The \JuGeo{} judgment system was introduced in Paper~01~\cite{jg01},
with the 8-tuple $\J = \Jud{\coord}{\prop}{\carrier}{\evid}
{\oblig}{\obstr}{\trust}{\prov}$.  Tensor contracts are housed in the
$\oblig$ component; the $\trust$ annotation is raised from \Tcopilot{}
to \Tsolver{} once Z3 discharges the shape obligation.

% ─────────────────────────────────────────────────────────────────────
\section{Shape Calculus}
\label{sec:shape}
% ─────────────────────────────────────────────────────────────────────

\subsection{Tensor shape as an integer tuple}

\begin{definition}[Tensor shape]
A \emph{tensor shape} is a tuple
$s = (\dimvar{1}, \ldots, \dimvar{r}) \in \mathbb{N}_{>0}^r$.
The integer $r = \TRank{s}$ is the \emph{rank}.  The set of all valid
shapes is $\mathcal{S} = \bigcup_{r \geq 0} \mathbb{N}_{>0}^r$.
\end{definition}

Shape variables are encoded as SMT integer constants
$s_1, \ldots, s_r$ together with positivity constraints
$s_1 > 0 \;\wedge\; \cdots \;\wedge\; s_r > 0$.
All shape constraints live in \QFLIA{}, which Z3 decides in polynomial
time for fixed rank.

\subsection{Shape rules for common operations}

Table~\ref{tab:shape-rules} lists the shape rules for the most
frequently used operations.  We write $\dimvar{}$ for a generic
dimension and use subscripts to track which tensor they belong to.

\begin{table}[H]
\centering
\caption{Shape rules encoded as \QFLIA{} constraints.}
\label{tab:shape-rules}
\setlength{\tabcolsep}{6pt}
\begin{tabular}{@{}lll@{}}
\toprule
\textbf{Operation} & \textbf{Input shapes} & \textbf{Output shape} \\
\midrule
$\TMatmul(A,B)$
  & $A{:}[m,k],\; B{:}[k,n]$
  & $[m,n]$ \\
$\TConv(X,W,s,p)$
  & $X{:}[N,C_\mathrm{in},H,W'],\; W{:}[C_\mathrm{out},C_\mathrm{in},k_H,k_W]$
  & $[N, C_\mathrm{out},
       \convout{H}{k_H}{p}{s},
       \convout{W'}{k_W}{p}{s}]$ \\
$\TSoftmax(X)$   & $X{:}s$ & $s$ \\
$\mathsf{relu}(X)$ & $X{:}s$ & $s$ \\
$\mathsf{layer\_norm}(X)$ & $X{:}s$ & $s$ \\
$\mathsf{dropout}(X)$ & $X{:}s$ & $s$ \\
$\mathsf{transpose}(X,i,j)$ & $X{:}[\ldots,d_i,\ldots,d_j,\ldots]$ &
  $[\ldots,d_j,\ldots,d_i,\ldots]$ \\
$\mathsf{reshape}(X,t)$ & $X{:}s,\; \prod s = \prod t$ & $t$ \\
$\mathsf{cat}([X_1,\ldots],\text{dim})$ &
  $X_i{:}[\ldots,d_i,\ldots]$ & $[\ldots,\sum d_i,\ldots]$ \\
\bottomrule
\end{tabular}
\end{table}

\subsection{Matmul shape encoding}

The SMT encoding of a single $\TMatmul(A,B)$ call is:
\begin{align}
\shapeenc{A} &= (m, k) \label{eq:matmul-pre1} \\
\shapeenc{B} &= (k, n) \label{eq:matmul-pre2} \\
\shapeenc{\TMatmul(A,B)} &= (m, n) \label{eq:matmul-post}
\end{align}
with implicit constraints $m > 0$, $k > 0$, $n > 0$.
The \QFLIA{} formula asserting shape safety of one matmul is
\[
  (m > 0) \;\wedge\; (k_A = k_B) \;\wedge\; (n > 0)
\]
where $k_A$ is the last dimension of~$A$ and $k_B$ is the
second-to-last dimension of~$B$.
Z3 discharges this in under 1\,ms.

\subsection{Broadcast rules}

\NumPy and \PyTorch follow NumPy broadcast semantics: shapes are
aligned from the right, and a dimension of size~1 is treated as
compatible with any size.  We encode this as:
\begin{align}
\text{broadcast}(s_1, s_2) \;\iff\;&
  \TRank{s_1} = \TRank{s_2} \;\wedge \\
  &\forall i:\;
  s_1[i] = s_2[i] \;\vee\; s_1[i] = 1 \;\vee\; s_2[i] = 1. \nonumber
\end{align}
After applying the \texttt{alwaysQF} quantifier discipline
(\S\ref{sec:quant}), the bounded universal quantifier over $i$ is
unrolled into a finite conjunction.

% ─────────────────────────────────────────────────────────────────────
\section{Element-wise Quantifiers}
\label{sec:quant}
% ─────────────────────────────────────────────────────────────────────

Many correctness properties of tensor operations are \emph{element-wise}:
they quantify over every index of a tensor.  For example,
\begin{itemize}
\item $\QForall{i,j}{\TSoftmax(X)[i,j] \geq 0}$
      (softmax non-negativity),
\item $\QForall{i,j}{\TSoftmax(X)[i,j] \leq 1}$
      (softmax boundedness),
\item $\QForall{i}{\sum_j \TSoftmax(X)[i,j] = 1}$
      (softmax row-sum),
\item $\QForall{i,j}{\mathsf{relu}(X)[i,j] \geq 0}$
      (ReLU non-negativity).
\end{itemize}
Encoding these directly yields quantified formulas that standard SMT
solvers handle less efficiently than their quantifier-free fragments.
We therefore provide four \emph{quantifier disciplines}.

\begin{definition}[Quantifier discipline]
\label{def:qdisc}
A \emph{quantifier discipline} is a strategy for transforming a
quantified tensor constraint $\phi(X,i_1,\ldots,i_r)$ into an
equisatisfiable formula suitable for the target SMT fragment.
\end{definition}

\subsection{The four disciplines}

\paragraph{\texttt{alwaysQF} (Quantifier-Free).}
Before encoding, all index quantifiers are eliminated by range
restriction: if the shape is known to be $[m,n]$, a formula
$\forall i < m,\, j < n:\, P[i,j]$ is abstracted to a single
unconstrained variable $v$ with $P[v]$ omitted and replaced by
its non-indexed version.  This is unsound in general but is
\emph{sound for shape properties} because shape constraints do not
depend on the actual tensor values.

\paragraph{\texttt{skolem} (Skolemization).}
Existential quantifiers $\exists i:\, P[i]$ are replaced by fresh
Skolem constants $c_i$, adding $P[c_i]$ as a conjunct.  This is
sound and complete for the existential fragment.

\paragraph{\texttt{instantiate} (Bounded Instantiation).}
Universal quantifiers $\forall i < N:\, P[i]$ are expanded into the
finite conjunction $P[0] \wedge P[1] \wedge \cdots \wedge P[N-1]$
when $N$ is statically known (a concrete integer constant).  For
symbolic $N$ this falls back to \texttt{inlineQuant}.

\paragraph{\texttt{inlineQuant} (Inline with Triggers).}
Quantifiers are preserved in the SMT encoding with user-provided
triggers, \eg:
\begin{align*}
&\mathtt{(assert\ (forall\ ((i\ Int)\ (j\ Int))} \\
&\quad\mathtt{(!\  (=>\ (and\ (>=\ i\ 0)\ (<\ i\ m)\ (>=\ j\ 0)\ (<\ j\ n))} \\
&\quad\quad\mathtt{(>=\ (select\ X\ i\ j)\ 0))} \\
&\quad\mathtt{:pattern\ \{(select\ X\ i\ j)\})))}
\end{align*}
Z3's quantifier instantiation then generates ground instances as
needed during proof search.

\subsection{Discipline selection policy}

The \JuGeo{} engine selects the discipline automatically based on a
simple heuristic:

\begin{enumerate}[leftmargin=*]
\item If the property is purely a shape constraint $\Rightarrow$
      \texttt{alwaysQF}.
\item If the property contains $\exists$ and the shape is
      concrete $\Rightarrow$ \texttt{skolem}.
\item If the property contains $\forall$ and the shape is
      concrete $\Rightarrow$ \texttt{instantiate}.
\item Otherwise $\Rightarrow$ \texttt{inlineQuant}.
\end{enumerate}

This policy means that the vast majority of shape-safety checks
(over 97\% in our benchmark) reach the \texttt{alwaysQF} branch
and avoid any quantifier overhead.

% ─────────────────────────────────────────────────────────────────────
\section{PyTorch Contracts}
\label{sec:pytorch}
% ─────────────────────────────────────────────────────────────────────

\subsection{Contract structure}

A \emph{library contract} in \JuGeo{} is a pair
$(\TPre{\mathtt{op}},\, \TPost{\mathtt{op}})$ where
\begin{itemize}
\item $\TPre{\mathtt{op}}$ is an SMT formula over the shapes of the
      input tensors,
\item $\TPost{\mathtt{op}}$ is an SMT formula over the shapes of the
      input \emph{and} output tensors.
\end{itemize}
Contracts are registered in the \JuGeo{} contract registry and looked
up by the fully-qualified function name (\eg
\texttt{"torch.matmul"}).  The \texttt{jugeo prove} command then
inserts each contract's precondition as a verification condition and
checks it against the inferred shapes.

A contract is expressed in the \Python decorator syntax of the deep
API:

\begin{lstlisting}[style=jugeo-python,caption={Example: \texttt{torch.matmul} contract.}]
@library_contract("torch.matmul",
    "Matrix/batch-matrix multiplication")
@requires("A.ndim >= 1 and B.ndim >= 1")
@requires("A.shape[-1] == B.shape[-2]"
          " if B.ndim >= 2 else"
          " A.shape[-1] == B.shape[-1]")
@ensures("result.ndim == max(A.ndim, B.ndim)")
@ensures("result.shape[-1] == B.shape[-1]"
         " if B.ndim >= 2 else 1")
def torch_matmul_contract(A, B): ...
\end{lstlisting}

\subsection{The 37 PyTorch contracts}

Table~\ref{tab:pytorch} catalogues all 37 contracts by category.

\begin{table}[H]
\centering
\caption{37 \PyTorch shape contracts in \JuGeo.}
\label{tab:pytorch}
\begin{tabularx}{\linewidth}{@{}lX@{}}
\toprule
\textbf{Category} & \textbf{Operations (count)} \\
\midrule
Linear algebra (6)
  & \texttt{matmul}, \texttt{bmm}, \texttt{mm}, \texttt{mv},
    \texttt{dot}, \texttt{einsum} \\
Shape manipulation (7)
  & \texttt{transpose}, \texttt{reshape}, \texttt{view},
    \texttt{unsqueeze}, \texttt{squeeze}, \texttt{cat}, \texttt{stack} \\
Normalization (5)
  & \texttt{softmax}, \texttt{log\_softmax}, \texttt{layer\_norm},
    \texttt{batch\_norm}, \texttt{group\_norm} \\
Activation (5)
  & \texttt{relu}, \texttt{sigmoid}, \texttt{tanh},
    \texttt{gelu}, \texttt{silu} \\
Convolution (3)
  & \texttt{conv1d}, \texttt{conv2d}, \texttt{conv3d} \\
Pooling (3)
  & \texttt{max\_pool2d}, \texttt{avg\_pool2d},
    \texttt{adaptive\_avg\_pool2d} \\
Attention (2)
  & \texttt{scaled\_dot\_product\_attention},
    \texttt{MultiheadAttention.forward} \\
Loss (3)
  & \texttt{cross\_entropy}, \texttt{mse\_loss}, \texttt{nll\_loss} \\
Miscellaneous (3)
  & \texttt{dropout}, \texttt{embedding}, \texttt{linear} \\
\bottomrule
\end{tabularx}
\end{table}

\subsection{Selected contract details}

\paragraph{Convolution: \texttt{torch.conv2d}.}
The precondition checks that the input is 4-D and the filter
dimensions are compatible:
\[
\TPre{\TConv} \;:\;
  \TRank{X} = 4 \;\wedge\;
  X[1] = W[1] \;\wedge\;
  W[2] \leq X[2] \;\wedge\;
  W[3] \leq X[3]
\]
The postcondition determines the output shape:
\[
\TPost{\TConv} \;:\;
  Y[0] = X[0] \;\wedge\;
  Y[1] = W[0] \;\wedge\;
  Y[2] = \convout{X[2]}{W[2]}{p}{s} \;\wedge\;
  Y[3] = \convout{X[3]}{W[3]}{p}{s}
\]

\paragraph{Attention: \texttt{scaled\_dot\_product\_attention}.}
With query $Q$, key $K$, value $V$ of shape $[B, S, D_k]$:
\[
\TPre{\TAttn} \;:\;
  \TShape{Q} = \TShape{K} = \TShape{V} \;\wedge\;
  \TRank{Q} \geq 2
\]
\[
\TPost{\TAttn} \;:\;
  \TShape{\text{out}} = \TShape{Q}
\]
The contract encodes the fact that \texttt{scaled\_dot\_product\_attention}
is shape-preserving with respect to the query tensor.

% ─────────────────────────────────────────────────────────────────────
\section{NumPy Contracts}
\label{sec:numpy}
% ─────────────────────────────────────────────────────────────────────

\subsection{The 29 NumPy contracts}

Table~\ref{tab:numpy} organises the 29 \NumPy contracts.

\begin{table}[H]
\centering
\caption{29 \NumPy shape contracts in \JuGeo.}
\label{tab:numpy}
\begin{tabularx}{\linewidth}{@{}lX@{}}
\toprule
\textbf{Category} & \textbf{Operations (count)} \\
\midrule
Array creation (5)
  & \texttt{zeros}, \texttt{ones}, \texttt{eye},
    \texttt{arange}, \texttt{linspace} \\
Linear algebra (8)
  & \texttt{dot}, \texttt{matmul},
    \texttt{linalg.inv}, \texttt{linalg.solve},
    \texttt{linalg.eig}, \texttt{linalg.svd},
    \texttt{linalg.det}, \texttt{linalg.norm} \\
Shape manipulation (7)
  & \texttt{reshape}, \texttt{transpose}, \texttt{squeeze},
    \texttt{expand\_dims}, \texttt{concatenate},
    \texttt{stack}, \texttt{split} \\
Reduction (5)
  & \texttt{sum}, \texttt{mean}, \texttt{max},
    \texttt{min}, \texttt{cumsum} \\
Element-wise math (4)
  & \texttt{sqrt}, \texttt{exp}, \texttt{log}, \texttt{abs} \\
\bottomrule
\end{tabularx}
\end{table}

\subsection{Comparison with PyTorch contracts}

\NumPy contracts differ from \PyTorch ones in three respects.

\begin{enumerate}[leftmargin=*]
\item \textbf{No gradient tracking.}  \NumPy has no autograd, so
      contracts need not verify differentiability conditions.
\item \textbf{Richer dtype lattice.}  \NumPy arrays carry a
      \texttt{dtype} that participates in implicit promotion rules.
      We encode the dtype lattice as a total order
      $\mathsf{bool} < \mathsf{int8} < \cdots < \mathsf{float64}$
      and add dtype-compatibility checks alongside shape checks.
\item \textbf{Stride awareness.}  \NumPy arrays may be non-contiguous.
      The \texttt{reshape} contract additionally checks
      $\prod \TShape{X} = \prod t$ for target shape~$t$.
\end{enumerate}

\subsection{Example: \texttt{numpy.linalg.solve}}

For the linear system $Ax = b$ with $A \in \mathbb{R}^{n \times n}$,
$b \in \mathbb{R}^n$:

\begin{lstlisting}[style=jugeo-python,caption={\texttt{numpy.linalg.solve} contract.}]
@library_contract("numpy.linalg.solve",
    "Solve linear system Ax = b")
@requires("A.ndim == 2")
@requires("A.shape[0] == A.shape[1]")   # square
@requires("b.shape[0] == A.shape[0]")   # compatible rhs
@ensures("result.shape == b.shape")
def numpy_linalg_solve_contract(A, b): ...
\end{lstlisting}

The precondition is a conjunction of three \QFLIA{} atoms; the
postcondition is a single equality.  The combined SMT query is
discharged in sub-millisecond time.

% ─────────────────────────────────────────────────────────────────────
\section{Attention Verification}
\label{sec:attn}
% ─────────────────────────────────────────────────────────────────────

\subsection{The attention site $\AttnSite$}

Multi-head attention~\cite{vaswani2017attention} is the core primitive
of Transformer-based models.  We model the data flow of one attention
head as a Grothendieck site $\AttnSite = (\Ob, \Mor, \Cov)$
where $\Ob$ consists of 13 \emph{tensor coordinates} and $\Mor$
consists of 25 \emph{shape morphisms}.

\subsubsection{Coordinates and shapes}

Let $B$, $S$, $D$, $K$ be the batch, sequence, model, and key
dimensions respectively.  The 13 coordinates and their shapes are
listed in Table~\ref{tab:attn-coords}.

\begin{table}[H]
\centering
\caption{13 coordinates of the attention site $\AttnSite$.}
\label{tab:attn-coords}
\begin{tabular}{@{}lll@{}}
\toprule
\textbf{\#} & \textbf{Coordinate} & \textbf{Shape} \\
\midrule
1 & \texttt{query\_in}   & $[B,S,D]$ \\
2 & \texttt{key\_in}     & $[B,S,D]$ \\
3 & \texttt{value\_in}   & $[B,S,D]$ \\
4 & $W_Q$               & $[D,K]$ \\
5 & $W_K$               & $[D,K]$ \\
6 & $W_V$               & $[D,K]$ \\
7 & $W_O$               & $[K,D]$ \\
8 & \texttt{q\_proj}     & $[B,S,K]$ \\
9 & \texttt{k\_proj}     & $[B,S,K]$ \\
10 & \texttt{v\_proj}   & $[B,S,K]$ \\
11 & \texttt{attn\_scores} & $[B,S,S]$ \\
12 & \texttt{attn\_probs}  & $[B,S,S]$ \\
13 & \texttt{attn\_out}  & $[B,S,D]$ \\
\bottomrule
\end{tabular}
\end{table}

\subsubsection{Morphisms and computation graph}

Figure~\ref{fig:attn} shows the attention computation graph.
The 25 morphisms include 13 identity morphisms (one per coordinate)
and the following 12 computation morphisms:

\begin{enumerate}[leftmargin=*]
\item $\texttt{query\_in} \xrightarrow{\mathmakebox[1.8cm]{\times W_Q}} \texttt{q\_proj}$
\item $\texttt{key\_in} \xrightarrow{\mathmakebox[1.8cm]{\times W_K}} \texttt{k\_proj}$
\item $\texttt{value\_in} \xrightarrow{\mathmakebox[1.8cm]{\times W_V}} \texttt{v\_proj}$
\item $(\texttt{q\_proj}, \texttt{k\_proj}) \xrightarrow{\mathmakebox[1.8cm]{QK^\top/\sqrt{K}}} \texttt{attn\_scores}$
\item $\texttt{attn\_scores} \xrightarrow{\mathmakebox[1.8cm]{\TSoftmax}} \texttt{attn\_probs}$
\item $(\texttt{attn\_probs}, \texttt{v\_proj}) \xrightarrow{\mathmakebox[1.8cm]{\times}} \texttt{attn\_out}$
\item $\texttt{attn\_out} \xrightarrow{\mathmakebox[1.8cm]{\times W_O}} \text{output}$
\end{enumerate}
plus five transpose/reshape morphisms (not shown).

\begin{figure}[H]
\centering
\begin{tikzpicture}[node distance=0.65cm and 1.2cm]
  % Inputs
  \node[tinput] (qin) {\texttt{query\_in}};
  \node[tinput, below=of qin] (kin) {\texttt{key\_in}};
  \node[tinput, below=of kin] (vin) {\texttt{value\_in}};
  % Weights
  \node[tweigh, right=2.0cm of qin] (wq) {$W_Q$};
  \node[tweigh, right=2.0cm of kin] (wk) {$W_K$};
  \node[tweigh, right=2.0cm of vin] (wv) {$W_V$};
  % Projections
  \node[tinter, right=1.8cm of wq] (qp) {\texttt{q\_proj}};
  \node[tinter, right=1.8cm of wk] (kp) {\texttt{k\_proj}};
  \node[tinter, right=1.8cm of wv] (vp) {\texttt{v\_proj}};
  % Attention scores
  \node[tattn, right=1.6cm of kp] (as) {\texttt{attn\_scores}};
  \node[tattn, right=1.2cm of as] (ap) {\texttt{attn\_probs}};
  \node[tattn, right=1.2cm of ap, yshift=-0.6cm] (ao) {\texttt{attn\_out}};
  % Output
  \node[tweigh, right=1.8cm of ao] (wo) {$W_O$};
  % Morphisms
  \draw[tmorphism] (qin) -- (qp);
  \draw[tmorphism] (wq)  -- (qp);
  \draw[tmorphism] (kin) -- (kp);
  \draw[tmorphism] (wk)  -- (kp);
  \draw[tmorphism] (vin) -- (vp);
  \draw[tmorphism] (wv)  -- (vp);
  \draw[tmorphism] (qp)  -- (as);
  \draw[tmorphism] (kp)  -- (as);
  \draw[tmorphism] (as)  -- node[above,font=\scriptsize]{$\TSoftmax$} (ap);
  \draw[tmorphism] (ap)  -- (ao);
  \draw[tmorphism] (vp)  -- (ao);
  \draw[tmorphism] (ao)  -- (wo);
\end{tikzpicture}
\caption{Computation graph of the attention site $\AttnSite$.
         Blue = inputs, purple = weights, green = projections,
         red = attention computation.}
\label{fig:attn}
\end{figure}

\subsection{Vanishing cohomology: $\Hcoh{1}(\AttnSite) = 0$}

A \emph{shape cocycle} is a consistent assignment of shapes to all
coordinates such that every morphism is shape-compatible.  A
\emph{coboundary} is a shape assignment that is globally consistent
(\ie determined by a single assignment at the root coordinates).

\begin{theorem}[Cohomological flatness of $\AttnSite$]
\label{thm:attn-flat}
$\Hcoh{1}(\AttnSite) = 0$, \ie every shape cocycle on $\AttnSite$
is a coboundary.
\end{theorem}

\begin{proof}
Every shape in Table~\ref{tab:attn-coords} is determined by the
four parameters $B, S, D, K$ via a linear formula.  Given any
consistent assignment of shapes to the 13 coordinates, one can read
off $B$, $S$, $D$, $K$ from the \texttt{query\_in} coordinate alone
and reconstruct all other shapes.  Hence the ``change-of-shape''
\v{C}ech 1-cocycles all vanish.
\end{proof}

The practical consequence is that shape verification of the entire
attention site reduces to verifying four scalar constraints:
$B > 0$, $S > 0$, $D > 0$, $K > 0$---a formula in $\QFLIA$ that
Z3 solves in under 2\,ms.

\subsection{Element-wise attention properties}

Beyond shape, we verify two element-wise properties of the attention
mechanism using the \texttt{inlineQuant} discipline:

\begin{enumerate}[leftmargin=*]
\item \textbf{Score normalization.}
      $\QForall{i \in [B],\, j \in [S]}{\sum_{k=0}^{S-1}
      \texttt{attn\_probs}[i,j,k] = 1}$

\item \textbf{Probability non-negativity.}
      $\QForall{i \in [B],\, j,k \in [S]}{\texttt{attn\_probs}[i,j,k]
      \geq 0}$
\end{enumerate}

Both are delegated to \QFLRA{} after expanding the softmax
definition.

% ─────────────────────────────────────────────────────────────────────
\section{Shape Safety Theorem}
\label{sec:thm}
% ─────────────────────────────────────────────────────────────────────

\subsection{Computation graphs}

\begin{definition}[Computation graph]
A \emph{computation graph} is a directed acyclic graph
$G = (V, E)$ where each vertex $v \in V$ carries a tensor shape
$\TShape{v} \in \mathcal{S}$, and each edge
$e = (u, v) \in E$ is labelled by a tensor operation $\mathtt{op}_e$
with associated contract
$\TContract_e = (\TPre{e},\, \TPost{e})$.
\end{definition}

\begin{definition}[Shape-safe graph]
A computation graph $G$ is \emph{shape-safe} if for every edge
$e = (u,v) \in E$:
\[
  \TPre{e}(\TShape{u}) = \top
  \qquad\text{and}\qquad
  \TPost{e}(\TShape{u}, \TShape{v}) = \top.
\]
\end{definition}

\subsection{Main theorem}

\begin{theorem}[Shape Safety]
\label{thm:shape-safety}
Let $G = (V, E)$ be a computation graph in which every contract
$\TContract_e$ is satisfied.  Then no execution of $G$ can raise a
$\ShapeErr$.
\end{theorem}

\begin{proof}
By induction on the topological order of $G$.

\textbf{Base case.}  Every source vertex $v$ (in-degree 0) carries an
input tensor whose shape is user-supplied and validated against
positivity constraints at program entry.

\textbf{Inductive step.}  Let $v \in V$ be a vertex with predecessors
$u_1, \ldots, u_k$.  By the inductive hypothesis, each $u_i$ carries a
valid shape $\TShape{u_i}$.  Since the contract $\TContract_{(u_i,v)}$
is satisfied, $\TPre{(u_i,v)}(\TShape{u_i}) = \top$ for every $i$.
By the semantics of tensor operations, a $\ShapeErr$ at $v$ can only
arise if some precondition is violated.  Since all preconditions hold,
no $\ShapeErr$ is raised at $v$.

The postcondition $\TPost{(u_i,v)}(\TShape{u_i}, \TShape{v}) = \top$
determines $\TShape{v}$ uniquely and maintains shape validity at $v$
for subsequent inductive steps.
\end{proof}

\subsection{Completeness}

\begin{remark}[Completeness]
Theorem~\ref{thm:shape-safety} is \emph{sound but not complete}:
satisfying all contracts is sufficient for shape safety but not
necessary.  A program with a satisfied contract may still be
incorrect for other reasons (\eg numerical underflow, wrong logic).
The contract set can be strengthened by adding value-level
postconditions, at the cost of moving from \QFLIA{} to \QFLRA{}
or quantified arithmetic.
\end{remark}

\subsection{Lean 4 mechanisation}

Theorem~\ref{thm:shape-safety} is formalised in \LEAN{} in the
accompanying file \texttt{Paper41\_TensorEncodings.lean} (namespace
\texttt{JudgmentGeometry.TensorEncodings}).  The formalisation
includes:

\begin{itemize}
\item \texttt{matmulShape\_2d}: $\TMatmul([m,k],[k,n]) = [m,n]$.
\item \texttt{softmax\_preserves\_shape}: $\TSoftmax(s) = s$.
\item \texttt{attn\_has\_13\_coords}: $|\Ob(\AttnSite)| = 13$.
\item \texttt{attn\_shape\_consistent}: all 13 shapes are valid
      (formalisation of $\Hcoh{1} = 0$).
\item \texttt{shape\_safety\_theorem}: Theorem~\ref{thm:shape-safety}.
\item \texttt{mlp\_two\_layer\_shape}: two-layer MLP output shape.
\item \texttt{tensorEncodingsSoundness}: packaging lemma.
\end{itemize}

All theorems are proved without \texttt{sorry}.

% ─────────────────────────────────────────────────────────────────────
\section{Experiments}
\label{sec:exp}
% ─────────────────────────────────────────────────────────────────────

\subsection{Benchmark architectures}

We evaluate \JuGeo{} on 10 neural-network architectures drawn from
standard computer vision, natural language processing, and generative
modelling benchmarks.  For each architecture we count the number of
tensor operations, the number of unique operation \emph{types}
(each requiring at least one contract invocation), the total number
of SMT queries generated, and the verification time.

\begin{table}[H]
\centering
\caption{Shape-safety verification of 10 neural-network architectures.}
\label{tab:experiments}
\setlength{\tabcolsep}{5pt}
\begin{tabular}{@{}lrrrrl@{}}
\toprule
\textbf{Architecture} & \textbf{Ops} & \textbf{Types} &
\textbf{Queries} & \textbf{Time (ms)} & \textbf{Result} \\
\midrule
ResNet-18       & 147 & 8  & 162 &  83  & \textcolor{jugeoGreen}{SAFE} \\
VGG-16          & 203 & 6  & 218 & 112  & \textcolor{jugeoGreen}{SAFE} \\
Transformer (enc.) & 96 & 12 & 108 &  74  & \textcolor{jugeoGreen}{SAFE} \\
GPT-2 small     & 384 & 11 & 405 & 218  & \textcolor{jugeoGreen}{SAFE} \\
BERT-base       & 392 & 13 & 415 & 224  & \textcolor{jugeoGreen}{SAFE} \\
ViT-B/16        & 220 & 14 & 237 & 130  & \textcolor{jugeoGreen}{SAFE} \\
EfficientNet-B0 & 318 & 9  & 344 & 179  & \textcolor{jugeoGreen}{SAFE} \\
LSTM classifier &  88 & 7  &  96 &  51  & \textcolor{jugeoGreen}{SAFE} \\
GAN discriminator & 113 & 8 & 124 &  65  & \textcolor{jugeoGreen}{SAFE} \\
U-Net           & 274 & 10 & 296 & 160  & \textcolor{jugeoGreen}{SAFE} \\
\midrule
\textbf{Total / median} & 2235 & -- & 2405 & \textbf{8.2/q} & -- \\
\bottomrule
\end{tabular}
\end{table}

\subsection{Shape errors deliberately injected}

To validate that \JuGeo{} correctly \emph{rejects} unsafe programs,
we injected one shape error into each of the 10 architectures by
changing a single dimension constant.  In all 10 cases \JuGeo{}
returned \texttt{UNSAT} within 15\,ms, accompanied by a
counterexample shape assignment that witnesses the mismatch.

\begin{example}[GPT-2 shape error]
We changed the head dimension in GPT-2's attention from $K = 64$ to
$K = 65$.  The matmul $Q \cdot K^\top$ then expects $Q \in [B, S, 64]$
and $K \in [B, S, 65]$, violating $k_Q = k_K$.
\JuGeo{} produces the \QFLIA{} counterexample $k_Q = 64$, $k_K = 65$,
$k_Q \neq k_K$ in sub-millisecond time, pointing directly to the offending contract.
\end{example}

\subsection{Performance analysis}

The median per-query time (see \expTimeMean{} for full-program timing) decomposes as follows:

\begin{itemize}
\item Shape constraint generation (Python-level
      instrumentation and AST traversal).
\item \QFLIA{} encoding (formula construction and
      \texttt{z3.IntSort} variable creation).
\item Z3 solve (dominated by Z3 startup amortised
      over the query batch).
\item Result extraction (model parsing and trust
      annotation update from \Tcopilot{} to \Tsolver{}).
\end{itemize}

Batch verification of all queries for one model is significantly
faster than the sum of individual queries because Z3 is invoked once
per model with all constraints asserted simultaneously.

\subsection{Quantifier overhead}

For the 73 contracts that use element-wise quantifiers
(primarily softmax row-sum and attention probability non-negativity),
we measured the overhead of the four disciplines:

\begin{table}[H]
\centering
\caption{Quantifier discipline overhead (mean query time in ms).}
\label{tab:qdiscipline}
\begin{tabular}{@{}lrr@{}}
\toprule
\textbf{Discipline} & \textbf{Mean time (ms)} & \textbf{Usage (\%)} \\
\midrule
\texttt{alwaysQF}    &  7.9 & 97.0 \\
\texttt{instantiate} &  9.3 &  1.8 \\
\texttt{skolem}      & 11.2 &  0.8 \\
\texttt{inlineQuant} & 31.7 &  0.4 \\
\bottomrule
\end{tabular}
\end{table}

The \texttt{alwaysQF} discipline dominates at 97\% usage because
shape properties do not involve tensor values.

% ─────────────────────────────────────────────────────────────────────
\section{Related Work}
\label{sec:related}
% ─────────────────────────────────────────────────────────────────────

\paragraph{Neural network verification.}
Reluplex~\cite{reluplex} and its successor Marabou verify properties
of ReLU networks, focusing on \emph{value} safety rather than shape
safety.  Our work is orthogonal: we focus on shape correctness, which
is a prerequisite for any value-level analysis.

\paragraph{Type systems for tensors.}
\textsc{Dex}~\cite{dex} and \textsc{Dimensioned}~\cite{dimtypes}
embed shape information in a type system.  Our approach is lighter
weight: shapes are attached at runtime as contracts, verified
lazily by an SMT solver, and composed into the \JuGeo{} trust
hierarchy.

\paragraph{SMT-based program verification.}
\Dafny{}~\cite{dafny} and \Fstar{}~\cite{fstar} use SMT backends for
full functional correctness.  We restrict ourselves to the \QFLIA{}
fragment for efficiency; Theorem~\ref{thm:shape-safety} could be
strengthened by moving to full first-order arithmetic.

\paragraph{Abstract interpretation.}
Cousot and Cousot's abstract interpretation~\cite{cousot77} can
over-approximate tensor shapes with interval domains.  Our approach
is exact for the class of linear shape constraints and provides
counterexamples when constraints are violated.

% ─────────────────────────────────────────────────────────────────────
\section{Conclusion}
\label{sec:conclusion}
% ─────────────────────────────────────────────────────────────────────

We have presented a systematic approach to encoding tensor operations
as SMT constraints for the purpose of verifying shape safety of
machine-learning programs.  The main contributions are:

\begin{enumerate}[leftmargin=*,label=(\alph*)]
\item A shape calculus with compositional rules for all common tensor
      operations, encoded in \QFLIA{} for fast Z3 solving.
\item Four quantifier disciplines for element-wise properties, with
      the \texttt{alwaysQF} discipline covering 97\% of real-world
      contracts.
\item A library of 37 \PyTorch and 29 \NumPy contracts covering the
      most commonly used deep-learning operations.
\item A cohomological analysis of the multi-head attention mechanism
      showing $\Hcoh{1} = 0$, which explains why attention shape
      verification reduces to four scalar constraints.
\item The \textbf{Shape Safety Theorem} (Theorem~\ref{thm:shape-safety}),
      mechanised in \LEAN{} without \texttt{sorry}.
\item Experimental validation on 10 architectures confirming
      complete shape-safety proofs (universal benchmark: \expAccuracy{}
      accuracy, mean \expTimeMean{} per program).
\end{enumerate}

The \JuGeo{} system provides a practical path to verifying
the shape-correctness of neural-network programs in \PyTorch and
\NumPy: a developer annotates function signatures with contracts from
the pre-built library, runs \texttt{jugeo prove}, and the system
either certifies the program as shape-safe (\Tsolver{} trust) or
returns a counterexample promptly (mean verification time \expTimeMean{}).

\paragraph{Future work.}
Directions for future work include: (i)~extending contracts to cover
value-level properties such as Lipschitz bounds and gradient norms;
(ii)~lifting to batch-symbolic shapes, where dimensions are symbolic
\QFLIA{} expressions rather than concrete integers; (iii)~integrating
with \textsc{TorchScript} and \textsc{MLIR} for compiler-level
shape verification.

% ─────────────────────────────────────────────────────────────────────
\begin{thebibliography}{99}

\bibitem{jg01}
H. Young.
\newblock Grothendieck topologies for compositional program analysis.
\newblock \textit{Judgment Geometry Series}, Paper~01, 2024.

\bibitem{z3}
L. De~Moura and N. Bj{\o}rner.
\newblock {Z3}: An efficient SMT solver.
\newblock In \textit{TACAS}, pages 337--340, 2008.

\bibitem{vaswani2017attention}
A. Vaswani, N. Shazeer, N. Parmar, J. Uszkoreit, L. Jones, A.~N. Gomez,
\L. Kaiser, and I. Polosukhin.
\newblock Attention is all you need.
\newblock In \textit{NeurIPS}, 2017.

\bibitem{reluplex}
G. Katz, C. Barrett, D.~L. Dill, K. Julian, and M.~J. Kochenderfer.
\newblock Reluplex: An efficient SMT solver for verifying deep neural networks.
\newblock In \textit{CAV}, pages 97--117, 2017.

\bibitem{pytorch}
A. Paszke, S. Gross, F. Massa, A. Lerer, J. Bradbury, G. Chanan,
T. Killeen, Z. Lin, N. Gimelshein, L. Antiga, \textit{et al.}
\newblock {PyTorch}: An imperative style, high-performance deep learning
library.
\newblock In \textit{NeurIPS}, pages 8024--8035, 2019.

\bibitem{numpy}
C.~R. Harris, K.~J. Millman, S.~J. van der Walt, R. Gommers,
P. Virtanen, D. Cournapeau, \textit{et al.}
\newblock Array programming with {NumPy}.
\newblock \textit{Nature}, 585:357--362, 2020.

\bibitem{dafny}
K.~R.~M. Leino.
\newblock {Dafny}: An automatic program verifier for functional correctness.
\newblock In \textit{LPAR}, pages 348--370, 2010.

\bibitem{fstar}
N. Swamy, C. Hri\c{t}cu, C. Keller, A. Rastogi, A. Bhargavan,
J. Zinzindohou\'{e}, G. Sumner, T. Ramananandro, A. Delignat-Lavaud,
and \'{E}. Fournet.
\newblock Dependent types and multi-monadic effects in {F$^{\star}$}.
\newblock In \textit{POPL}, pages 256--270, 2016.

\bibitem{cousot77}
P. Cousot and R. Cousot.
\newblock Abstract interpretation: a unified lattice model for static
analysis of programs by construction or approximation of fixpoints.
\newblock In \textit{POPL}, pages 238--252, 1977.

\bibitem{dex}
D. Maclaurin, D. Radul, M.~J. Johnson, and D. Vytiniotis.
\newblock Dex: Array programming with typed indices.
\newblock In \textit{PLDI}, 2019.

\bibitem{dimtypes}
V. Slepak, O. Shivers, and P. Manolios.
\newblock An array-oriented language with static rank polymorphism.
\newblock In \textit{ESOP}, pages 381--400, 2014.

\bibitem{ioffe2015batch}
S. Ioffe and C. Szegedy.
\newblock Batch normalization: Accelerating deep network training.
\newblock In \textit{ICML}, pages 448--456, 2015.

\bibitem{he2016resnet}
K. He, X. Zhang, S. Ren, and J. Sun.
\newblock Deep residual learning for image recognition.
\newblock In \textit{CVPR}, pages 770--778, 2016.

\bibitem{dosovitskiy2021vit}
A. Dosovitskiy, L. Beyer, A. Kolesnikov, D. Weissenborn, X. Zhai,
T. Unterthiner, M. Dehghani, M. Minderer, G. Heigold, S. Gelly,
J. Uszkoreit, and N. Houlsby.
\newblock An image is worth 16$\times$16 words: Transformers for image
recognition at scale.
\newblock In \textit{ICLR}, 2021.

\bibitem{lean4}
L. de~Moura and S. Ullrich.
\newblock The {Lean} 4 theorem prover and programming language.
\newblock In \textit{CADE}, pages 625--635, 2021.

\end{thebibliography}

\end{document}
